\chapter{数字と時間 — Números y Tiempo}
\unidadheader{02}{数字と時間}{Números y Tiempo}{Hablar de horas, fechas y precios}

\begin{tcolorbox}[vocabbox, title={📌 Vocabulario Nuevo}]
\begin{tabularx}{\linewidth}{llX}
\toprule
\textbf{Japonés} & \textbf{Español} & \textbf{Tipo} \\
\midrule
\vocabitem{〜時}{じ}{hora ~ (ej: 三時 = 3:00)}{contador}
\vocabitem{〜分}{ふん・ぷん}{minuto ~ (ej: 五分 = 5 min)}{contador}
\vocabitem{午前}{ごぜん}{AM / mañana}{sust.}
\vocabitem{午後}{ごご}{PM / tarde}{sust.}
\vocabitem{半}{はん}{y media (30 min)}{sust.}
\vocabitem{今}{いま}{ahora}{adv.}
\vocabitem{〜年}{ねん}{año ~ (ej: 二〇二六年)}{contador}
\vocabitem{〜月}{がつ}{mes ~ (ej: 五月 = mayo)}{contador}
\vocabitem{〜日}{にち・か}{día ~ (ej: 三日 = día 3)}{contador}
\vocabitem{〜曜日}{ようび}{día de semana (月・火・水・木・金・土・日)}{sust.}
\bottomrule
\end{tabularx}
\end{tcolorbox}

\begin{tcolorbox}[phrasebox, title={🔢 Números 1–10000}]
\begin{tabular}{llllll}
一 いち & 二 に & 三 さん & 四 し/よん & 五 ご & 六 ろく \\
七 なな/しち & 八 はち & 九 く/きゅう & 十 じゅう & 百 ひゃく & 千 せん \\
\end{tabular}\\[4pt]
\small 百 (100) · 千 (1000) · 万 (10000) · 十万 (100000)
\end{tcolorbox}

\subsection*{🗣️ Frases Clave}

\frase{\ruby{今}{いま}\ruby{何時}{なんじ}ですか？}{¿Qué hora es ahora?}
\frase{\ruby{午後}{ごご}\ruby{三時}{さんじ}\ruby{半}{はん}です。}{Son las 3:30 PM.}
\frase{\ruby{今日}{きょう}は\ruby{何月何日}{なんがつなんにち}ですか？}
  {¿A qué fecha estamos hoy?}
\frase{\ruby{誕生日}{たんじょうび}はいつですか？}{¿Cuándo es tu cumpleaños?}

\subsection*{⚙️ Gramática}

\patron{Expresar hora}
  {〜時〜分です}
  {\ej{\ruby{午前}{ごぜん}\ruby{七時}{しちじ}\ruby{十五分}{じゅうごふん}です。}{Son las 7:15 AM.}
  \ej{\ruby{十二時}{じゅうにじ}\ruby{半}{はん}です。}{Son las 12:30.}}

\patron{Expresar fecha}
  {〜年〜月〜日（〜曜日）}
  {\ej{二〇二六\ruby{年}{ねん}\ruby{五月}{ごがつ}\ruby{十八日}{じゅうはちにち}\ruby{月曜日}{げつようび}です。}
  {Es lunes, 18 de mayo de 2026.}}

\patron{Desde... hasta}
  {〜から〜まで}
  {\ej{\ruby{九時}{くじ}から\ruby{五時}{ごじ}まではたらきます。}
  {Trabajo de 9 a 5.}}

\begin{tcolorbox}[ejerciciobox, title={📝 Ejercicios Rápidos}]
\begin{enumerate}
  \item ¿Cómo se dice "Son las 8:45 PM"?
  \item Escribe tu fecha de cumpleaños en japonés.
  \item ¿Cuánto es 三千八百円?
\end{enumerate}
\vspace{0.4em}
\textbf{Respuestas:}
1) \ruby{午後}{ごご}\ruby{八時}{はちじ}\ruby{四十五分}{よんじゅうごふん}です。\quad
2) (personal)\quad
3) 3800 yen.
\end{tcolorbox}

\begin{tcolorbox}[kanjibox, title={🀄 Kanjis de la Unidad}]
\begin{tabularx}{\linewidth}{cllXX}
\toprule
\textbf{Kanji} & \textbf{On} & \textbf{Kun} & \textbf{Significado} & \textbf{Vocab} \\
\midrule
\kanjirow{時}{じ}{とき}{tiempo / hora}{\ruby{時間}{じかん} / \ruby{何時}{なんじ}}
\kanjirow{分}{ふん・ぶん}{わ}{minuto / dividir}{\ruby{五分}{ごふん} / \ruby{半分}{はんぶん}}
\kanjirow{年}{ねん}{とし}{año}{\ruby{今年}{ことし} / \ruby{来年}{らいねん}}
\kanjirow{月}{がつ・つき}{つき}{mes / luna}{\ruby{一月}{いちがつ} / \ruby{月曜日}{げつようび}}
\kanjirow{日}{にち・じつ}{ひ・か}{día / sol}{\ruby{今日}{きょう} / \ruby{誕生日}{たんじょうび}}
\kanjirow{今}{こん・きん}{いま}{ahora / este}{\ruby{今日}{きょう} / \ruby{今}{いま}}
\bottomrule
\end{tabularx}
\end{tcolorbox}
